\documentclass[10pt,a4paper]{amsart}
\usepackage{amsmath,amssymb,amsthm,graphicx}
\usepackage{verbatim}
\newenvironment{reference}{\comment}{\endcomment}
%\newenvironment{reference}{}{}
\newenvironment{slogan}{\comment}{\endcomment}
\newenvironment{history}{\comment}{\endcomment}

% For commutative diagrams we use Xy-pic
\usepackage[all]{xy}

% We use 2cell for 2-commutative diagrams.
\xyoption{2cell}
% Theorem environments.
%
\theoremstyle{plain}
\newtheorem{theorem}[subsection]{Theorem}
\newtheorem{proposition}[subsection]{Proposition}
\newtheorem{lemma}[subsection]{Lemma}

\theoremstyle{definition}
\newtheorem{definition}[subsection]{Definition}
\newtheorem{example}[subsection]{Example}
\newtheorem{exercise}[subsection]{Exercise}
\newtheorem{situation}[subsection]{Situation}

\theoremstyle{remark}
\newtheorem{remark}[subsection]{Remark}
\newtheorem{remarks}[subsection]{Remarks}

\numberwithin{equation}{subsection}

% Macros
%
\def\lim{\mathop{\mathrm{lim}}\nolimits}
\def\colim{\mathop{\mathrm{colim}}\nolimits}
\def\Spec{\mathop{\mathrm{Spec}}}
\def\Hom{\mathop{\mathrm{Hom}}\nolimits}
\def\Ext{\mathop{\mathrm{Ext}}\nolimits}
\def\SheafHom{\mathop{\mathcal{H}\!\mathit{om}}\nolimits}
\def\SheafExt{\mathop{\mathcal{E}\!\mathit{xt}}\nolimits}
\def\Sch{\mathit{Sch}}
\def\Mor{\mathop{\mathrm{Mor}}\nolimits}
\def\Ob{\mathop{\mathrm{Ob}}\nolimits}
\def\Sh{\mathop{\mathit{Sh}}\nolimits}
\def\NL{\mathop{N\!L}\nolimits}
\def\CH{\mathop{\mathrm{CH}}\nolimits}
\def\proetale{{pro\text{-}\acute{e}tale}}
\def\etale{{\acute{e}tale}}
\def\QCoh{\mathit{QCoh}}
\def\Ker{\mathop{\mathrm{Ker}}}
\def\Im{\mathop{\mathrm{Im}}}
\def\Coker{\mathop{\mathrm{Coker}}}
\def\Coim{\mathop{\mathrm{Coim}}}

% Boxtimes
%
\DeclareMathSymbol{\boxtimes}{\mathbin}{AMSa}{"02}

%
% Macros for moduli stacks/spaces
%
\def\QCohstack{\mathcal{QC}\!\mathit{oh}}
\def\Cohstack{\mathcal{C}\!\mathit{oh}}
\def\Spacesstack{\mathcal{S}\!\mathit{paces}}
\def\Quotfunctor{\mathrm{Quot}}
\def\Hilbfunctor{\mathrm{Hilb}}
\def\Curvesstack{\mathcal{C}\!\mathit{urves}}
\def\Polarizedstack{\mathcal{P}\!\mathit{olarized}}
\def\Complexesstack{\mathcal{C}\!\mathit{omplexes}}
% \Pic is the operator that assigns to X its picard group, usage \Pic(X)
% \Picardstack_{X/B} denotes the Picard stack of X over B
% \Picardfunctor_{X/B} denotes the Picard functor of X over B
\def\Pic{\mathop{\mathrm{Pic}}\nolimits}
\def\Picardstack{\mathcal{P}\!\mathit{ic}}
\def\Picardfunctor{\mathrm{Pic}}
\def\Deformationcategory{\mathcal{D}\!\mathit{ef}}
\usepackage{cleveref}
\crefname{originalreference}{original reference}{original references}
\Crefname{originalreference}{Original reference}{Original references}
\makeatletter
\newlabel{algebra-proposition-Auslander-Buchsbaum}{{090V}{source}{}{sourcepointer-1}{}}
\newlabel{algebra-proposition-Auslander-Buchsbaum@cref}{{[originalreference][0][]090V}{[1][0][]source}{}{}{}}
\makeatother
\crefname{originalreference}{original reference}{original references}
\Crefname{originalreference}{Original reference}{Original references}
\makeatletter
\newlabel{algebra-lemma-acyclic}{{00N0}{source}{}{sourcepointer-2}{}}
\newlabel{algebra-lemma-acyclic@cref}{{[originalreference][0][]00N0}{[1][0][]source}{}{}{}}
\makeatother
\title{CM intersection multiplicities equal lengths of scheme-theoretic intersections}
\author{The Stacks Project Authors}
\date{stacks-project:intersection.tex:local-snapshot-20261012; GNU FDL 1.2 or later}
\begin{document}
\maketitle
\noindent\textit{Editorial note.} All delivered mathematical text is the author's own, copied verbatim from the pinned Stacks Project chapter file intersection.tex; no statement or proof step was rewritten, reordered or removed.\par\medskip
\noindent\textit{Editorial note.} Span recomputation by label pairing (LANE\_BRIEF 10.2): in this chapter the author's proof environment follows the statement environment as a standalone block; the spec covers the complete environment pair (statement span, proof span) as recorded in the pool candidate.\par\medskip
\noindent\textit{Editorial note.} Occupation scan: combined claim\_key + statement\_anchor + generic-line scan over all registration files (active and withdrawn) found 0 hits for this chapter (is-occupied.json, 2026-10-14); consistent with the duty manager's six-check dispatch.\par\medskip
\ifdefined\hypertarget\hypertarget{sourcepointer-1}{}\fi
\noindent\textit{Source reference.} Original reference 090V: sources/stacks-chow/tags.txt, source lines 11423--11423. Cross-chapter standard prerequisite cited by the author inside the counted span; it is used as a pointer and its own proof is not part of this item.\par\medskip
\ifdefined\hypertarget\hypertarget{sourcepointer-2}{}\fi
\noindent\textit{Source reference.} Original reference 00N0: sources/stacks-chow/tags.txt, source lines 842--842. Cross-chapter standard prerequisite cited by the author inside the counted span; it is used as a pointer and its own proof is not part of this item.\par\medskip
\begin{lemma}
\label{lemma-intersection-multiplicity-CM}
Let $X$ be a nonsingular variety and $W, V \subset X$ closed
subvarieties which intersect properly. Let $Z$ be an irreducible component
of $V \cap W$ with generic point $\xi$. Assume that $\mathcal{O}_{W, \xi}$
and $\mathcal{O}_{V, \xi}$ are Cohen-Macaulay. Then
$$
e(X, V \cdot W, Z) =
\text{length}_{\mathcal{O}_{X, \xi}}(\mathcal{O}_{V \cap W, \xi})
$$
where $V \cap W$ is the scheme theoretic intersection.
In particular, if both $V$ and $W$ are Cohen-Macaulay, then
$V \cdot W = [V \cap W]_{\dim(V) + \dim(W) - \dim(X)}$.
\end{lemma}

\begin{proof}
Set $A = \mathcal{O}_{X, \xi}$, $B = \mathcal{O}_{V, \xi}$, and
$C = \mathcal{O}_{W, \xi}$. By Auslander-Buchsbaum
(Algebra, Proposition \ref{algebra-proposition-Auslander-Buchsbaum})
we can find a finite free resolution $F_\bullet \to B$ of
length
$$
\text{depth}(A) - \text{depth}(B) =
\dim(A) - \dim(B) = \dim(C)
$$
First equality as $A$ and $B$ are Cohen-Macaulay and the second
as $V$ and $W$ intersect properly. Then $F_\bullet \otimes_A C$ is a
complex of finite free modules representing $B \otimes_A^\mathbf{L} C$
hence has cohomology modules with support in $\{\mathfrak m_A\}$.
By the Acyclicity lemma (Algebra, Lemma \ref{algebra-lemma-acyclic})
which applies as $C$ is Cohen-Macaulay
we conclude that $F_\bullet \otimes_A C$ has nonzero
cohomology only in degree $0$. This finishes the proof.
\end{proof}

\end{document}
